(1) 弦的正向入射与反射位移分别 \(y=e^{i(kx-\omega t)}+re^{i(-kx-\omega t)}\)，其中 \(k=\omega/c\)。右端弦对端点的张力沿 x 方向的横向合力是 \(-T y_x(0,t)\)，减振器力为 \(-b y_t(0,t)\)。端点无质量，故 \(Ty_x+b y_t=0\)。以复幅值代入得到 \(iTk(1-r)-i b\omega(1+r)=0\)。

(2) 令 \(Z=Tk/\omega=\sqrt{T\mu}\)，解得
\[
r=\frac{Z-b}{Z+b}.
\]
反射率 \(\mathcal R=|r|^2=(Z-b)^2/(Z+b)^2\)，吸收率为
\[
\mathcal A=1-\mathcal R=\frac{4Zb}{(Z+b)^2}.
\]
以波的入射平均功率与端点 \(b\overline{y_t^2}\) 比较，也可直接得到同一吸收率。

(3) 当 b=Z，边界与行波特征机械阻抗完全匹配，r=0、\(\mathcal A=1\)，无反射。对 b>0，恒有 \((Z-b)^2\ge0\)，所以 \(\mathcal A\le1\) 且 b=Z 时取等号。

(4) b→0 为自由端 r→+1，位移无相位反转；b→∞ 为固定端 r→−1，位移反相。两极限均 \(\mathcal A\to0\)，与端点无法消耗平均机械功一致。
